% Figure: a hypodermic syringe as a Poiseuille-flow device. A plunger pushed at
% force F drives fluid from a wide barrel of radius R_b through a long thin
% needle of radius R_n and length L_n. Almost the entire pressure drop falls
% across the needle, because the Poiseuille resistance scales as R^{-4} and the
% needle bore is far smaller than the barrel bore. Drawn as a schematic
% cross-section with the barrel, plunger, and needle labelled.
\begin{figure}[htbp]
\centering
\begin{tikzpicture}[thick]
% Barrel walls (a wide rectangle).
\draw[engblack] (0,-0.9) rectangle (5,0.9);
% Fluid fill in the barrel.
\fill[engblue!14] (1.6,-0.85) rectangle (5,0.85);
% Plunger (a shaded block) with a push rod.
\fill[enggray!30] (1.0,-0.85) rectangle (1.6,0.85);
\draw[engblack] (1.0,-0.85) rectangle (1.6,0.85);
\draw[engblack, line width=1.4pt] (0.2,0) -- (1.0,0);
\draw[->, engred, line width=1.2pt] (-0.9,0) -- (0.15,0);
\node[font=\scriptsize, engred, anchor=south] at (-0.5,0.06) {$F$};
% Barrel radius label.
\draw[<->, enggray] (5.15,-0.9) -- (5.15,0.9);
\node[font=\scriptsize, enggray, anchor=west] at (5.2,0) {$R_b$};
% Needle: a long thin channel to the right.
\draw[engblack] (5,0.12) -- (9.5,0.12);
\draw[engblack] (5,-0.12) -- (9.5,-0.12);
\fill[engblue!14] (5,-0.12) rectangle (9.5,0.12);
% Needle bore and length labels.
\draw[<->, enggray] (9.7,-0.12) -- (9.7,0.12);
\node[font=\scriptsize, enggray, anchor=west] at (9.75,0) {$R_n$};
\draw[<->, enggray] (5,-0.5) -- (9.5,-0.5);
\node[font=\scriptsize, enggray, anchor=north] at (7.25,-0.55) {$L_n$ (needle)};
% Exit flow arrow.
\draw[->, engblue, line width=1.1pt] (9.5,0) -- (10.4,0);
\node[font=\scriptsize, engblue, anchor=south] at (10.0,0.06) {$Q$};
% Pressure-drop annotation across the needle.
\node[font=\scriptsize, engred, anchor=south] at (7.25,0.35)
  {$\Delta p \approx \dfrac{8\mu L_n Q}{\pi R_n^4}$};
\end{tikzpicture}
\caption[The hypodermic syringe as a Poiseuille device]{A hypodermic syringe
read as a Hagen-Poiseuille resistance. The plunger force $F$ drives a flow $Q$
from the wide barrel of radius $R_b$ through the long thin needle of radius
$R_n$ and length $L_n$. Because the Poiseuille resistance scales as the inverse
fourth power of the bore, almost the whole pressure drop falls across the
needle, whose bore is far smaller than the barrel's, and the barrel contributes
almost none. The plunger force needed is this pressure drop times the barrel
area, so a small needle demands a large force for a modest flow. Worked
example~35 carries the numbers for an insulin injection.}
\label{fig:vol03ch12:syringe}
\end{figure}
